% =========================================================
% RESUMEN / ABSTRACT
% =========================================================
\chapter*{Resumen}
\addcontentsline{toc}{chapter}{Resumen}
Contexto: La formación de planetas terrestres mediante acreción de pebbles depende críticamente del transporte radial de sólidos en discos protoplanetarios. Sin embargo, la presencia de subestructuras observadas en estos discos puede modificar dicho transporte, alterando tanto el crecimiento planetario como la entrega de agua hacia las regiones internas.

Objetivo: En este trabajo se investiga cómo las trampas de presión y la migración temporal de la línea de nieve regulan el flujo de pebbles y afectan la formación de planetas terrestres y de waterworlds.

Metodología: Para ello se desarrolló un framework auto-consistente basado en TriPoDPy para modelar la evolución de gas y polvo, junto con un módulo de pebble accretion (PA3Py) para calcular el crecimiento de embriones planetarios sobre miles de simulaciones que exploran diferentes configuraciones de turbulencia, subestructuras y propiedades del polvo.

Resultados: Se identificó una \textit{Isla de Crecimiento} en el espacio 
de parámetros ($\alpha \lesssim 10^{-3}$, $M_{\rm gap} \sim 0.1\,M_{\rm Jup}$) 
donde ocurre formación planetaria eficiente. El mecanismo clave es la 
\textit{regulación temporal del flujo de pebbles}: las subestructuras 
controlan no solo cuánto material llega al planeta, sino cuándo lo hace.

Esta modulación temporal genera una correlación entre masa y contenido 
de agua mediante el pebble snow amplificado: liberación tardía de 
reservorios de polvo helado cuando la snowline migra hacia el interior. 
Para muy baja turbulencia ($\alpha = 10^{-4}$), emerge una bifurcación 
poblacional: planetas secos vs ricos en agua, dependiendo de la 
cronología relativa del crecimiento.

Conclusiones: Las subestructuras regulan simultáneamente crecimiento 
y composición planetaria mediante modulación temporal del flujo de 
pebbles. La cronología es tan importante como el inventario de sólidos. 
Estos resultados explican la formación eficiente de planetas terrestres 
ricos en agua en discos estructurados, con predicciones testables 
mediante ALMA.